% !TeX root = ../main.tex
\chapter{Methodology and Experimental Pipeline}

\section{Overview of the Experimental Pipeline}

The core of this work is a systematic experimental pipeline designed to compare various neural evaluation architectures and state-space partitioning strategies under the constraints of embedded chess engines. Rather than focusing on a single novel algorithm in isolation, this thesis evaluates a battery of techniques that take inspiration from other fields, like LLMs and robotics. They range from standard baselines, like vanilla NNUE, to advanced Mixture-of-Experts (MoE) formulations. The pipeline is structured around a rigorous loop of architectural design, hyperparameter selection, large-scale training, and precise evaluation against a strong teacher signal.

The foundation of the entire pipeline is a large-scale dataset of approximately 150 million chess positions. These positions are extracted from human games and labeled with high-quality Win-Draw-Loss (WDL) probability distributions by the Leela Chess Zero (Lc0) engine at search depth 1. This provides a rich, calibrated soft-target distribution that captures the teacher's uncertainty, serving as the ground truth for all subsequent training and evaluation \parencite{lczero}.

The experimental methodology follows three distinct phases for every model variant:

\begin{enumerate}
\item \textbf{Hyperparameter Selection and Scaling:} Determining the optimal architectural widths, learning rates, and data budgets. This includes running dataset scaling tests to understand when the base architectures saturate, which directly informs the data allocation for MoE expert fine-tuning.
\item \textbf{Training:} Executing the training procedures, which vary depending on the method (e.g., joint end-to-end training for soft MoE, versus a multi-stage pipeline of base training, clustering, and expert fine-tuning for hard MoE).
\item \textbf{Evaluation:} Assessing the trained models primarily through the \textbf{Cross-Entropy (CE)} against the teacher's WDL distribution on a strictly held-out test set. This metric directly quantifies how well the student model replicates the teacher's probabilistic understanding of the position, serving as the definitive criterion for model selection and comparison.
\end{enumerate}

\section{Base Evaluation Architectures}

Before exploring complex routing and partitioning strategies, it is necessary to establish strong, efficient baselines. The pipeline evaluates two distinct approaches to the core evaluation function.

\subsection{Standard NNUE}

The primary baseline is the Efficiently Updatable Neural Network (NNUE). This architecture relies on a sparse, handcrafted feature encoding (piece-square tables) mapped through a shared accumulator layer (L1), followed by a dense hidden layer (L2) and a linear output head. The dual-perspective design ensures the network always evaluates the board from the side-to-move's viewpoint. The NNUE baseline represents the state-of-the-art for resource-constrained engines, offering a highly favorable trade-off between evaluation strength and incremental update speed \parencite{nasu2018}.

\subsection{NNUE with a Learned L1 World Model}

To investigate whether the handcrafted sparse accumulator can be replaced by a learned representation, the pipeline also explores a variant where the conventional L1 accumulator is substituted with a learned "world model" encoder. Instead of relying on predefined piece-square features, this architecture utilizes a feedforward neural network encoder that processes the raw board state. The hypothesis is that shifting the learning effort entirely to a shared, learned representation might yield a more adaptable spatial understanding of the board, albeit at the cost of losing the incremental update efficiency of the standard NNUE accumulator.

\section{Experts and State-Space Partitioning}

To test the hypothesis that different regions of the state space require specialized evaluation heads, the pipeline extends the base architectures with Mixture-of-Experts (MoE) formulations. In all MoE variants, the total parameter budget is distributed across multiple expert heads, and a routing mechanism directs each position to the appropriate expert(s). The pipeline explores four distinct routing and partitioning strategies, ranging from simple heuristics to fully differentiable gating networks \parencite{shazeer2017}.

\subsection{Hard MoE with Piece-Count Clustering}

The most traditional approach to NNUE bucketing uses handcrafted, domain-specific heuristics. In this variant, the state space is partitioned into discrete buckets based on the scalar piece count (the total material remaining on the board). This serves as the conventional baseline for hard MoE, representing the standard practice in engines like Stockfish. The routing is deterministic, computationally free, and requires no learned dispatcher, but it relies on the assumption that game phase (as approximated by material) is the primary driver of evaluation divergence \parencite{stockfish}.

\subsection{Hard MoE with L1 Activation Clustering}

Moving away from handcrafted rules, this variant employs unsupervised clustering on the internal representations of the base model. By applying Mini-Batch K-Means to the frozen L1 accumulator activations, positions that "look similar" to the network's feature extractor are grouped together. A lightweight dispatcher is trained to predict these bucket assignments at inference time. This method tests whether partitioning based on the model's current knowledge (its representations) yields better specialization than handcrafted heuristics \parencite{sculley2010}.

\subsection{Hard MoE with Sample-Gradient Clustering}

% [Gradient Clustering]
To test whether the optimization dynamics of the model can guide expert specialization, this thesis explores a partitioning method based on per-sample gradients. Rather than clustering what the model currently \emph{knows} (its hidden activations), this approach clusters what the model \emph{needs to adjust} to better fit each individual position. By grouping positions that demand similar parameter updates, the resulting buckets are inherently aligned with the learning signal, theoretically minimizing optimization conflicts during expert fine-tuning. The pipeline for this variant consists of five stages:

\begin{enumerate}
	\item \textbf{Base Training:} A reference NNUE is trained to convergence on the full dataset. This establishes a stable parameter space and provides the frozen L1 representation required for subsequent routing.
	\item \textbf{Gradient Extraction:} For every position in the dataset, the per-sample gradient of the loss is computed with respect to the \emph{trainable head parameters} (excluding the frozen L1 accumulator). These vectors are L2-normalized to focus the clustering purely on the direction of the required update, discarding magnitude information which is sensitive to position difficulty.
	\item \textbf{Clustering:} Mini-Batch K-Means is applied to the normalized gradient vectors to discover $B$ distinct buckets of positions exhibiting aligned learning requirements.
	\item \textbf{Dispatcher Training:} Since computing gradients at inference time is computationally prohibitive and requires the unknown teacher evaluation, a lightweight dispatcher must be trained to predict the gradient bucket assignment using only the frozen L1 activations. This step introduces a fundamental geometric challenge, as the gradient decision boundaries are often non-convex and scattered within the linear L1 activation space.
	\item \textbf{Expert Fine-Tuning:} The base head is cloned $B$ times, and each expert is fine-tuned exclusively on the positions assigned to its corresponding gradient bucket, while the shared L1 accumulator remains frozen.
\end{enumerate}

\subsection{End-to-End Soft and Sparse MoE}

Finally, the pipeline evaluates fully differentiable routing mechanisms trained jointly with the experts. This includes \textbf{Soft-Gated MoE} (where the router outputs a softmax distribution and the top-$k$ experts are blended) and \textbf{Switch MoE} (a sparse, hard top-1 routing mechanism). Both are trained end-to-end with an auxiliary load-balancing loss to prevent expert collapse. This variant tests whether a jointly optimized routing policy can discover partitions superior to the offline, two-stage clustering methods, despite the strict computational constraints of the target hardware \parencite{shazeer2017} \parencite{fedus2021}.


\subsection{KITSUNE: Knowledge-guided Initialization and Two-stage Segregation}
\label{sec:kitsune}

The experimental results presented in Chapter 5 reveal a fundamental geometric limitation in single-stage routing: while sample gradients provide the most informative signal for expert specialization, they form non-convex, scattered manifolds in the activation space that are poorly approximated by linear dispatchers. Conversely, L1 activations form convex, linearly separable regions but lack direct alignment with the learning signal. To bridge this gap without sacrificing inference efficiency, we propose \textbf{KITSUNE} (\textbf{K}nowledge-guided \textbf{I}nitialization and \textbf{T}wo-stage \textbf{S}egregation for \textbf{U}nsupervised \textbf{N}eural \textbf{E}xperts).

KITSUNE decouples the geometric partitioning of the activation space from the semantic assignment of experts via a two-stage offline pipeline, reducing the online inference cost to a strictly linear operation plus a static lookup.

\subsubsection*{Stage 1: Micro-Quantization of L1 Space}
Instead of clustering directly into $B$ macro-experts, the L1 activation space is first partitioned into a large number of fine-grained \emph{micro-buckets} ($B'$), where $B' \gg B$ (e.g., $B'=64$ for $B=8$). Because these micro-buckets represent small, locally convex regions of the activation manifold, they can be recovered at inference time by a lightweight linear dispatcher with near-perfect accuracy ($>99\%$). This stage ensures that the online routing remains robust and computationally trivial.

\subsubsection*{Stage 2: Gradient-Informed Stitching}
The semantic specialization is introduced offline through a \emph{stitching} process. For each micro-bucket, we compute the mean sample gradient with respect to the head parameters. These mean gradients serve as high-fidelity proxies for the learning requirements of the positions within that micro-bucket. We then apply agglomerative clustering on these $B'$ gradient vectors to merge them into $B$ \emph{macro-experts}. 

This stitching step allows the formation of arbitrarily non-convex and disconnected routing regions in the L1 space, as a single macro-expert may be composed of multiple disparate micro-buckets that share similar gradient directions. To ensure computational balance, the stitching algorithm enforces a constraint that merges micro-buckets in multiples, preventing any single expert from becoming a bottleneck during parallel evaluation.

\subsubsection*{Inference and Benefits}
At inference time, the KITSUNE router operates in two steps: (1) the linear dispatcher predicts the micro-bucket index from the L1 activations, and (2) a static lookup table maps this index to the final macro-expert. This architecture retains the actionable gradient signal that drives expert specialization—addressing the failure mode of standard gradient clustering—while preserving the zero-overhead, deterministic routing of L1-based methods. Furthermore, by initializing experts with positions that are not only geometrically close but also gradient-aligned, KITSUNE provides a more coherent starting point for the fine-tuning phase, potentially accelerating convergence and improving final evaluation strength.


\section{Training and Evaluation Protocol}

To ensure a fair and rigorous comparison across all architectures and routing strategies, the pipeline enforces a strict training and evaluation protocol.

\subsection{Hyperparameter Selection and Dataset Scaling}

Before committing to the final MoE experiments, the pipeline includes a scaling phase to determine the data requirements of the base architectures. By training models of varying widths on progressively larger subsets of the 150M dataset, the pipeline identifies the point at which the base NNUE saturates. This scaling analysis is critical for the MoE variants: it reveals the exact data budget required per expert to avoid underfitting, and highlights the tension between the massive dataset available for base training and the fraction of data each expert receives in a hard MoE setup.

\subsection{Training Procedures}

Base models and end-to-end soft MoE variants are trained using the Adam optimizer \parencite{kingma2015} with a linearly decaying learning rate. The loss function is universally the \textbf{soft cross-entropy} between the model's predicted WDL distribution and the Lc0 teacher distribution. 

For the hard MoE variants, the training is split into a base phase (using Adam) and an expert fine-tuning phase. Crucially, during the expert fine-tuning phase, the optimizer is switched to \textbf{Rprop} (resilient backpropagation). As detailed in the experimental results, Rprop's sign-only, per-weight step-size adaptation proves significantly more stable than Adam for fine-tuning the small, specialized expert heads, preventing the premature convergence and instability observed with adaptive gradient methods.

\begin{itemize}
	\item For \textbf{base models and end-to-end MoE}, the training relies entirely on Adam, with the load-balancing loss weighted to ensure uniform expert utilization in the soft-gated variants.
	\item For \textbf{L1-clustered and gradient-clustered hard MoE}, the expert fine-tuning phase utilizes Rprop to ensure stable specialization, while the L1 accumulator remains frozen.
	\item The \textbf{piece-count} variant requires no routing training, as the buckets are fixed by definition.
\end{itemize}

\subsection{Evaluation Metrics}

The primary metric for model selection and comparison is the \textbf{Cross-Entropy (CE)} evaluated on a strictly held-out test set of positions. CE directly measures the divergence between the student's predicted WDL distribution and the teacher's true distribution, capturing both the accuracy and the calibrated uncertainty of the evaluation. Additionally, the pipeline tracks:

\begin{itemize}
	\item \textbf{Mean Absolute Error (MAE)} on the scalar expected value ($p_W - p_L$), which correlates with the actual search behavior of the engine.
	\item \textbf{Routing Diagnostics:} For the hard MoE variants, the pipeline evaluates the dispatcher's Top-1 accuracy, as well as clustering agreement metrics such as the Adjusted Rand Index \parencite{hubert1985} and Normalized Mutual Information \parencite{strehl2002}. These quantify how faithfully the deployable routing signal (L1 activations) recovers the offline partition (especially for the gradient-based method).
\end{itemize}
